\chapter{Fundamentação Teórica}
\label{chapter:theorical_found}

A crescente complexidade dos sistemas industriais~\cite{grieves2014digital} modernos tem impulsionado o desenvolvimento de \ac{dt} como ferramentas avançadas de monitoramento, análise e apoio à tomada de decisão. No contexto da indústria de óleo e gás, esses sistemas computacionais buscam representar, em ambiente virtual, o comportamento de ativos físicos e seus processos operacionais, permitindo a análise contínua do estado do sistema, a antecipação de falhas e a otimização das operações. Nesse cenário, modelos de aprendizado de máquina têm desempenhado um papel central na construção de gêmeos digitais, especialmente na modelagem de fenômenos dinâmicos complexos que não podem ser facilmente descritos por modelos físicos tradicionais.

Grande parte dos dados utilizados na construção de gêmeos digitais industriais é composta por séries temporais provenientes de sensores e sistemas de monitoramento, que registram continuamente variáveis operacionais dos ativos. Dessa forma, torna-se natural a aplicação de métodos de aprendizado de máquina capazes de capturar dependências temporais e padrões dinâmicos presentes nesses dados. Entre esses métodos, destacam-se as \ac{rnn} e suas variantes mais avançadas, que foram projetadas especificamente para lidar com dados sequenciais e dependências de longo prazo.

Neste contexto, este capítulo apresenta uma introdução técnica aos principais modelos de aprendizado de máquina que serão utilizados no desenvolvimento do gêmeo digital proposto neste projeto. Inicialmente, serão discutidas arquiteturas baseadas em redes neurais recorrentes, com destaque para os modelos \ac{lstm} e \ac{gru}, amplamente utilizados em tarefas de previsão e análise de séries temporais devido à sua capacidade de mitigar problemas clássicos de treinamento, como o vanishing gradient, e capturar relações temporais complexas em sequências de dados.

Além das arquiteturas recorrentes, o capítulo também aborda modelos baseados em mecanismos de atenção, incluindo a arquitetura Transformer, que tem se destacado como uma alternativa poderosa para o processamento de dados sequenciais. Esses modelos utilizam mecanismos de atenção para identificar relações relevantes entre diferentes elementos de uma sequência, permitindo capturar dependências de longo alcance de forma mais eficiente do que muitas arquiteturas recorrentes tradicionais.

Outro aspecto relevante do projeto está relacionado à distribuição geográfica dos ativos monitorados, característica comum em ambientes de produção da indústria de óleo e gás, especialmente em operações offshore. Nesses cenários, a centralização de dados pode apresentar limitações associadas a latência, largura de banda, requisitos de segurança da informação e restrições operacionais. Para lidar com essas limitações, este capítulo também apresenta os fundamentos do aprendizado federado (Federated Learning), uma abordagem de treinamento distribuído na qual múltiplos nós ou ativos contribuem para o treinamento de um modelo global sem a necessidade de compartilhamento direto dos dados brutos.

Assim, este capítulo tem como objetivo fornecer a base conceitual e técnica necessária para compreender as principais arquiteturas de aprendizado de máquina que serão utilizadas no desenvolvimento do gêmeo digital deste projeto. A discussão apresentada servirá como fundamento para as etapas subsequentes de modelagem, implementação e validação dos modelos preditivos aplicados ao monitoramento e à análise de ativos industriais no setor de óleo e gás.

\section{Aprendizado Federado}